\subsection{Divisores e ideales fraccionarios invertibles}
\label{subsection:IV.21.2}

\begin{env}[21.2.1]
\label{IV.21.2.1}
Sea $(X,\sh{O}_X)$ un espacio anillado.
Se llama \emph{ideal fraccionario sobre $X$} a un sub-$\sh{O}_X$-módulo del $\sh{O}_X$-módulo $\sh{M}_X$ de los gérmenes de funciones meromorfas sobre $X$.
Un ideal fraccionario $\sh{J}$ sobre $X$ que sea un $\sh{O}_X$-módulo invertible se llama \emph{ideal fraccionario invertible}.
\end{env}

\begin{proposition}[21.2.2]
\label{IV.21.2.2}
Para que un ideal fraccionario $\sh{J}$ sobre $X$ sea invertible, es necesario y suficiente que para todo $x\in X$ exista un entorno abierto $U$ de $x$ y una sección $f\in\Gamma(U,\sh{M}_X^*)$ tales que $\sh{J}|U=\sh{O}_U.f$.
\end{proposition}

\begin{proof}
La condición es evidentemente suficiente, pues la aplicación $s\mapsto s(f|V)$ de $\Gamma(V,\sh{O}_X)$ en $\Gamma(V,\sh{J})$ es evidentemente biyectiva para todo abierto $V\subset U$.
Para ver que es necesaria, observemos que por hipótesis existe un entorno abierto $U$ de $x$ y un isomorfismo de $\sh{O}_X$-módulos $\sh{O}_U\xrightarrow{\sim}\sh{J}|U$.
Si $f$ es la imagen de la sección $1\in\Gamma(U,\sh{O}_X)$ por este isomorfismo, se puede suponer, restringiendo $U$, que $f=u/s$, donde $u\in\Gamma(U,\sh{O}_X)$ y $s\in\Gamma(U,\sh{S}(\sh{O}_X))$, y el isomorfismo considerado hace corresponder entonces, a toda sección $v\in\Gamma(V,\sh{O}_X)$ (donde $V$ es un abierto contenido en $U$), la sección $v(u|V)/(s|V)$;
decir que la aplicación así definida es biyectiva significa que $u|V$ es un elemento regular de $\Gamma(V,\sh{O}_X)$, y por tanto $f\in\Gamma(U,\sh{M}_X^*)$.
\end{proof}

Obsérvese que para todo abierto $U$ de $X$ tal que $\sh{J}|U=\sh{O}_U.f$, con $f\in\Gamma(U,\sh{M}_X^*)$, la sección $f$ está determinada \emph{salvo multiplicación} por un elemento de $\Gamma(U,\sh{O}_X^*)$, puesto que la
\oldpage[IV-4]{259}
multiplicación por estos elementos proporciona todos los automorfismos del $\sh{O}_U$-módulo $\sh{O}_U$ (\EGAReadLink{0.5.4.3}{$\mathbf{0}_{\mathrm I}$, 5.4.3}).

\begin{corollary}[21.2.3]
\label{IV.21.2.3}
\begin{enumerate}[label=(\roman*)]
  \item Sea $\sh{J}$ un ideal fraccionario invertible;
  entonces el $\sh{O}_X$-módulo invertible $\sh{J}^{-1}$ se identifica canónicamente con el ideal fraccionario $\sh{J}'$ (\emph{transportador} de $\sh{J}$ en $\sh{O}_X$) definido del modo siguiente: para todo abierto $U$ de $X$ tal que $\sh{J}|U=\sh{O}_U.f$, donde $f\in\Gamma(U,\sh{M}_X^*)$, se tiene $\sh{J}'|U=\sh{O}_U.f^{-1}$.
  \item Si $\sh{J}_1$ y $\sh{J}_2$ son dos ideales fraccionarios invertibles, la aplicación canónica $\sh{J}_1\otimes\sh{J}_2\to\sh{J}_1\sh{J}_2$ es un isomorfismo de $\sh{O}_X$-módulos.
\end{enumerate}
\end{corollary}

\begin{proof}
La afirmación (ii) resulta inmediatamente de \sref{IV.21.2.2}.
Por otra parte, la observación hecha al final de \sref{IV.21.2.2} prueba que existe efectivamente un único ideal fraccionario $\sh{J}'$ definido por la condición del enunciado;
el isomorfismo canónico de $\sh{J}'$ sobre $\sh{J}^{-1}=\shHom_{\sh{O}_X}(\sh{J},\sh{O}_X)$ se obtiene haciendo corresponder a toda sección $s(f^{-1}|V)$ de $\Gamma(V,\sh{J}')$ (donde $V$ es un abierto contenido en $U$ y $s\in\Gamma(V,\sh{O}_X)$) el homomorfismo $t(f|V)\mapsto st$ de $\Gamma(V,\sh{J})$ en $\Gamma(V,\sh{O}_X)$.
\end{proof}

En virtud de \sref{IV.21.2.3},~(i), se identificarán generalmente los $\sh{O}_X$-módulos invertibles $\sh{J}'$ y $\sh{J}^{-1}$, considerando por tanto $\sh{J}^{-1}$ como un sub-$\sh{O}_X$-módulo de $\sh{M}_X$.

\begin{env}[21.2.4]
\label{IV.21.2.4}
De \sref{IV.21.2.3} resulta que el conjunto $\operatorname{Id.inv}(X)$ de los ideales fraccionarios invertibles sobre $X$ está dotado de una estructura de \emph{grupo conmutativo} para la ley de composición $(\sh{J}_1,\sh{J}_2)\mapsto\sh{J}_1\sh{J}_2$, siendo $\sh{O}_X$ el elemento neutro de este grupo.
Es claro que para todo abierto $U$ de $X$ se tiene $(\sh{J}_1\sh{J}_2)|U=(\sh{J}_1|U)(\sh{J}_2|U)$, por lo que la aplicación de restricción $\sh{J}\mapsto\sh{J}|U$ es un homomorfismo de grupos $\operatorname{Id.inv}(X)\to\operatorname{Id.inv}(U)$;
se define así un \emph{prehaz de grupos conmutativos} $U\mapsto\operatorname{Id.inv}(U)$;
es inmediato que, de hecho, este prehaz es un \emph{haz de grupos conmutativos}, que se denota $\sh{J}d.inv_X$.
\end{env}

\begin{env}[21.2.5]
\label{IV.21.2.5}
Para toda función meromorfa regular $f\in\Gamma(X,\sh{M}_X^*)$, de \sref{IV.21.2.2} resulta que $\sh{J}(f)=\sh{O}_X.f$ es un ideal fraccionario invertible, y se tiene evidentemente $\sh{J}(f_1f_2)=\sh{J}(f_1)\sh{J}(f_2)$; en otras palabras, la aplicación $f\mapsto\sh{J}(f)$ es un homomorfismo del grupo conmutativo $\Gamma(X,\sh{M}_X^*)$ en el grupo conmutativo $\operatorname{Id.inv}(X)$.
Sustituyendo $X$ por un abierto cualquiera $U$ de $X$ y observando que los homomorfismos obtenidos son compatibles con las operaciones de restricción, se obtiene un homomorfismo canónico de haces de grupos conmutativos:
\[
\label{IV.21.2.5.1}
  I_0:\sh{M}_X^*\to\sh{J}d.inv_X.
  \tag{21.2.5.1}
\]
Obsérvese que si $f\in\Gamma(X,\sh{O}_X^*)$, se tiene $\sh{J}(f)=\sh{O}_X$;
de ello se deduce inmediatamente que el homomorfismo $I_0$ se factoriza como
\[
\label{IV.21.2.5.2}
  \sh{M}_X^*\to\sh{M}_X^*/\sh{O}_X^*=\sh{D}iv_X\xrightarrow{I}\sh{J}d.inv_X
  \tag{21.2.5.2}
\]
donde $I$ es un homomorfismo del haz de grupos aditivos $\sh{D}iv_X$ en el haz de grupos multiplicativos $\sh{J}d.inv_X$;
se tiene, por tanto, para todo abierto $U$ de $X$ un homomorfismo $\sh{J}_U:\operatorname{Div}(U)\to\operatorname{Id.inv}(U)$ de grupos conmutativos tal que, para toda sección $f\in\Gamma(U,\sh{M}_X^*)$, se tiene
\[
\label{IV.21.2.5.3}
  \sh{J}_U(\operatorname{div}_U(f))=\sh{O}_U.f.
  \tag{21.2.5.3}
\]

\oldpage[IV-4]{260}

Se concluye que $\sh{J}_X(D)$, para todo divisor $D\in\operatorname{Div}(X)$, es el ideal fraccionario invertible definido del modo siguiente: para todo abierto $U$ de $X$ tal que $D|U=\operatorname{div}_U(f)$, donde $f\in\Gamma(U,\sh{M}_X^*)$, $\sh{J}_X(D)|U$ es el ideal fraccionario invertible $\sh{O}_U.f$.
Se tiene, por tanto, en virtud de \sref{IV.21.1.6}, para toda función meromorfa regular $f\in\Gamma(X,\sh{M}_X^*)$, la relación
\[
\label{IV.21.2.5.4}
  f\in\Gamma(X,\sh{J}_X(D)) \iff \operatorname{div}(f)\geq D.
  \tag{21.2.5.4}
\]
\end{env}

\begin{proposition}[21.2.6]
\label{IV.21.2.6}
El homomorfismo $I:\sh{D}iv_X\to\sh{J}d.inv_X$ es biyectivo.
\end{proposition}

\begin{proof}
Se define, en efecto, un homomorfismo $I_X'$ de $\operatorname{Id.inv}(X)$ en $\operatorname{Div}(X)$ haciendo corresponder a todo ideal fraccionario invertible $\sh{J}$ sobre $X$ el divisor $I_X'(\sh{J})$ definido del modo siguiente: para todo abierto $U$ de $X$ tal que $\sh{J}|U=\sh{O}_U.f$, donde $f\in\Gamma(U,\sh{M}_X^*)$ \sref{IV.21.2.2}, se toma $I_X'(\sh{J})|U=\operatorname{div}_U(f)$;
en virtud de la observación que sigue a \sref{IV.21.2.2}, esta definición es independiente del generador $f$ elegido en $\sh{J}|U$ y determina efectivamente un divisor sobre $X$.
Además, esta definición muestra inmediatamente que los homomorfismos $\sh{J}_X$ e $I_X'$ son inversos entre sí.
Sustituyendo $X$ por un abierto cualquiera $U$, se deduce la definición del isomorfismo $I':\sh{J}d.inv_X\to\sh{D}iv_X$ inverso de $I$, de donde resulta la proposición.
Se pondrá $I_X'(\sh{J})=\operatorname{div}(\sh{J})$, de modo que, para toda función meromorfa regular $f$ sobre $X$, se tiene
\[
\label{IV.21.2.6.1}
  \operatorname{div}(\sh{O}_X.f)=\operatorname{div}(f).
  \tag{21.2.6.1}
\]
\end{proof}

\begin{env}[21.2.7]
\label{IV.21.2.7}
Se identificarán a menudo los haces $\sh{D}iv_X$ y $\sh{J}d.inv_X$ (respectivamente, los grupos $\operatorname{Div}(X)$ y $\operatorname{Id.inv}(X)$) mediante los isomorfismos anteriores $I$ e $I'$ (respectivamente, $\sh{J}_X$ e $I_X'$).
Obsérvese que se tiene la relación
\[
\label{IV.21.2.7.1}
  D\geq0 \iff \sh{J}_X(D)\subset\sh{O}_X \quad\text{para}\quad D\in\operatorname{Div}(X)
  \tag{21.2.7.1}
\]
como resulta inmediatamente de las definiciones \sref{IV.21.1.6} y \sref{IV.21.1.5.1};
en otras palabras, la imagen $\sh{J}_X(\operatorname{Div}^+(X))$ es el conjunto de los \emph{ideales} de $\sh{O}_X$ (llamados también a veces \emph{ideales enteros}) que son $\sh{O}_X$-módulos invertibles: tal ideal $\sh{J}$ se caracteriza además por el hecho de que para todo $x\in X$ existe un entorno abierto $U$ de $x$ en $X$ tal que $\sh{J}|U=\sh{O}_U.f$, donde $f$ es un elemento regular de $\Gamma(U,\sh{O}_X)$.
El conjunto $\sh{J}_X(\operatorname{Div}^+(X))$ de estos ideales es, por tanto, un submonoide de $\operatorname{Id.inv}(X)$, igual al conjunto de los elementos positivos para una relación de orden compatible con la estructura de grupo de $\operatorname{Id.inv}(X)$, y es inmediato que esta relación no es otra que la relación \emph{opuesta} a la inclusión;
en otras palabras, se tiene
\[
\label{IV.21.2.7.2}
  \sh{J}_1\subset\sh{J}_2 \iff \operatorname{div}(\sh{J}_1)\geq\operatorname{div}(\sh{J}_2)
  \tag{21.2.7.2}
\]
para $\sh{J}_1$, $\sh{J}_2$ en $\operatorname{Id.inv}(X)$.
\end{env}

\begin{env}[21.2.8]
\label{IV.21.2.8}
Para todo divisor $D$ sobre $X$, se pone
\[
\label{IV.21.2.8.1}
  \sh{O}_X(D)=(\sh{J}_X(D))^{-1};
  \tag{21.2.8.1}
\]
$\sh{O}_X(D)$ es, por tanto, un ideal fraccionario invertible definido del modo siguiente: para todo abierto $U$ de $X$ tal que $D|U=\operatorname{div}_U(f)$, donde $f\in\Gamma(U,\sh{M}_X^*)$, $\sh{O}_X(D)|U$ es el ideal
\oldpage[IV-4]{261}
fraccionario invertible $\sh{O}_U.f^{-1}$;
en virtud de \sref{IV.21.1.6}, para toda función meromorfa regular $f$, se tiene la relación
\[
\label{IV.21.2.8.2}
  f\in\Gamma(X,\sh{O}_X(D)) \iff \operatorname{div}(f)\geq-D.
  \tag{21.2.8.2}
\]
Además, es claro que se tienen los isomorfismos canónicos \sref{IV.21.2.3}
\[
\label{IV.21.2.8.3}
  \left\{\begin{aligned}
    &\sh{O}_X(0)=\sh{O}_X,\quad \sh{O}_X(D+D')=\sh{O}_X(D)\sh{O}_X(D')\xrightarrow{\sim}\sh{O}_X(D)\otimes\sh{O}_X(D')\\
    &\sh{O}_X(nD)=(\sh{O}_X(D))^n\xrightarrow{\sim}(\sh{O}_X(D))^{\otimes n}
  \end{aligned}\right.
  \tag{21.2.8.3}
\]
para todo entero $n\in\mathbf{Z}$ y para dos divisores cualesquiera $D$, $D'$ sobre $X$.
\end{env}

\begin{env}[21.2.9]
\label{IV.21.2.9}
Sea $\sh{J}$ un ideal fraccionario invertible sobre $X$.
La inyección canónica $\sh{J}\hookrightarrow\sh{M}_X$ define por tensorización un homomorfismo de $\sh{O}_X$-módulos
\[
\label{IV.21.2.9.1}
  \sh{M}_X(\sh{J})=\sh{J}\otimes_{\sh{O}_X}\sh{M}_X\to\sh{M}_X\otimes_{\sh{O}_X}\sh{M}_X=\sh{M}_X
  \tag{21.2.9.1}
\]
que es un \emph{isomorfismo}: en efecto, si $U$ es un abierto de $X$ tal que $\sh{J}|U=\sh{O}_U.f$, donde $f\in\Gamma(U,\sh{M}_X^*)$, el homomorfismo \sref{IV.21.2.9.1} restringido a $U$ no es otro que el isomorfismo que a toda sección $t$ de $\sh{M}_X(\sh{J})|U=\sh{M}_X|U$ sobre $V\subset U$ hace corresponder la sección $t(f|V)^{-1}$ del mismo haz.
En el isomorfismo \sref{IV.21.2.9.1}, a las secciones meromorfas regulares de $\sh{J}$ sobre $X$ corresponden las funciones meromorfas regulares sobre $X$.

Consideremos en particular el caso en que $\sh{J}=\sh{O}_X(D)$, donde $D$ es un divisor sobre $X$;
se tiene entonces un isomorfismo canónico
\[
\label{IV.21.2.9.2}
  \sh{M}_X(\sh{O}_X(D))\xrightarrow{\sim}\sh{M}_X
  \tag{21.2.9.2}
\]
y se denota por $s_D$ la sección meromorfa regular de $\sh{O}_X(D)$ sobre $X$ que corresponde mediante este isomorfismo a la sección $1$ de $\sh{M}_X$.
Si $U$ es un abierto de $X$ tal que $\sh{O}_X(D)|U=\sh{O}_U.f^{-1}$, donde $f\in\Gamma(U,\sh{M}_X^*)$, se tiene $s_D|U=1$ en $\Gamma(U,\sh{M}_X)$;
como entonces se tiene $D|U=\operatorname{div}_U(f)$, se deduce de \sref{IV.21.1.4} que
\[
\label{IV.21.2.9.3}
  \operatorname{div}(s_D)=D.
  \tag{21.2.9.3}
\]
Por otra parte, de los isomorfismos canónicos \sref{IV.21.2.8.3} se deducen inmediatamente las fórmulas
\[
\label{IV.21.2.9.4}
  s_0=1,\quad s_{D+D'}=s_D\otimes s_{D'},\quad s_{nD}=s_D^{\otimes n}\quad(n\in\mathbf{Z}).
  \tag{21.2.9.4}
\]
\end{env}

\begin{env}[21.2.10]
\label{IV.21.2.10}
Consideremos, entre dos pares $(\sh{L},s)$, $(\sh{L}',s')$, donde $\sh{L}$ y $\sh{L}'$ son dos $\sh{O}_X$-módulos invertibles, y $s$ (respectivamente, $s'$) es una sección meromorfa regular de $\sh{L}$ (respectivamente, $\sh{L}'$) sobre $X$, la relación: «existe un isomorfismo $u:\sh{L}\xrightarrow{\sim}\sh{L}'$ tal que $\overline{u}(s)=s'$», donde $\overline{u}:\Gamma(X,\sh{M}_X(\sh{L}))\xrightarrow{\sim}\Gamma(X,\sh{M}_X(\sh{L}'))$ es el isomorfismo deducido de $u$ (obsérvese que el isomorfismo $u$ que satisface esta condición queda entonces determinado de manera única).
Es claro que se trata de una relación de equivalencia, y puesto que existe un conjunto de $\sh{O}_X$-módulos invertibles tal que todo $\sh{O}_X$-módulo invertible es isomorfo a un elemento de este conjunto (\EGAReadLink{0.5.4.7}{$\mathbf{0}_{\mathrm I}$, 5.4.7}), se puede hablar del conjunto $D(X)$ de las clases de equivalencia de los pares $(\sh{L},s)$ para la relación anterior.
Para todo $\sh{O}_X$-módulo invertible $\sh{L}$ y toda sección
\oldpage[IV-4]{262}
meromorfa regular $s$ de $\sh{L}$ sobre $X$, se denotará por $\operatorname{cl}(\sh{L},s)$ el elemento de $D(X)$ correspondiente al par $(\sh{L},s)$.
De (\EGAReadLink{0.5.4.3}{$\mathbf{0}_{\mathrm I}$, 5.4.3}) resulta que si $s$, $s'$ son dos secciones meromorfas regulares de $\sh{L}$ sobre $X$, la relación $\operatorname{cl}(\sh{L},s)=\operatorname{cl}(\sh{L},s')$ equivale a la existencia de una sección $t\in\Gamma(X,\sh{O}_X^*)$ tal que $s'=ts$.

Es inmediato que si $(\sh{L},s)$ es equivalente a $(\sh{L}_1,s_1)$ y $(\sh{L}',s')$ a $(\sh{L}_1',s_1')$, los pares $(\sh{L}\otimes\sh{L}',s\otimes s')$ y $(\sh{L}_1\otimes\sh{L}_1',s_1\otimes s_1')$ son equivalentes;
se define, por tanto, en $D(X)$ una ley de composición poniendo
\[
\label{IV.21.2.10.1}
  \operatorname{cl}(\sh{L},s)\operatorname{cl}(\sh{L}',s')=\operatorname{cl}(\sh{L}\otimes\sh{L}',s\otimes s');
  \tag{21.2.10.1}
\]
es inmediato que se trata de una ley de grupo conmutativo, cuyo elemento neutro es $\operatorname{cl}(\sh{O}_X,1)$ y en el que el inverso de $\operatorname{cl}(\sh{L},s)$ es $\operatorname{cl}(\sh{L}^{-1},s^{\otimes(-1)})$.
\end{env}

\begin{proposition}[21.2.11]
\label{IV.21.2.11}
Las aplicaciones
\[
\label{IV.21.2.11.1}
  D\mapsto\operatorname{cl}(\sh{O}_X(D),s_D),\qquad \operatorname{cl}(\sh{L},s)\mapsto\operatorname{div}(s)
  \tag{21.2.11.1}
\]
son isomorfismos inversos, respectivamente, de $\operatorname{Div}(X)$ sobre $D(X)$ y de $D(X)$ sobre $\operatorname{Div}(X)$.
\end{proposition}

\begin{proof}
Teniendo en cuenta \sref{IV.21.2.8.3}, \sref{IV.21.2.9.4} y \sref{IV.21.2.10.1}, basta ver que los compuestos de estas dos aplicaciones son las identidades en $\operatorname{Div}(X)$ y $D(X)$, respectivamente.
La primera afirmación no es otra que \sref{IV.21.2.9.3}.
Por otra parte, sea $D=\operatorname{div}(s)$, donde $s$ es una sección meromorfa regular de $\sh{L}$ sobre $X$, y sea $U$ un abierto de $X$ tal que exista un isomorfismo de $\sh{L}|U$ sobre $\sh{O}_U$ que transforme $s|U$ en $f\in\Gamma(U,\sh{M}_X^*)$, de modo que $D|U=\operatorname{div}_U(f)$, $\sh{O}_X(D)|U=\sh{O}_U.f^{-1}$ y $s_D|U$ es el elemento unidad de $\Gamma(U,\sh{M}_X)$.
Existe, por tanto, un isomorfismo $v_U:\sh{L}|U\to\sh{O}_U.f^{-1}=\sh{O}_X(D)|U$ tal que $\overline{v}_U$ (notación de \sref{IV.21.2.10}) transforma $s|U$ en $f.f^{-1}=1$;
se ve inmediatamente que estos isomorfismos son compatibles con las operaciones de restricción, por lo que definen un isomorfismo $v:\sh{L}\to\sh{O}_X(D)$ tal que $\overline{v}(s)=s_D$.
\end{proof}

Se puede transportar mediante el primero de los isomorfismos \sref{IV.21.2.11.1} la estructura de grupo ordenado de $\operatorname{Div}(X)$ a $D(X)$;
los elementos $\geq0$ de $D(X)$ son, por tanto, las clases $\operatorname{cl}(\sh{L},s)$ tales que $\operatorname{div}(s)\geq0$, es decir, \sref{IV.21.1.7.1}, tales que
\[
  s\in\Gamma(X,\sh{L})\cap\Gamma(X,(\sh{M}_X(\sh{L}))^*).
\]

\begin{env}[21.2.12]
\label{IV.21.2.12}
Sea $D$ un divisor \emph{positivo} sobre un \emph{preesquema} $X$;
el ideal fraccionario $\sh{J}_X(D)$ es entonces un ideal de $\sh{O}_X$ que es un $\sh{O}_X$-módulo invertible;
sea $Y(D)$ el subpreesquema cerrado de $X$ que este define.
Para todo $x\in Y(D)$ existe, por hipótesis, un entorno abierto $U$ de $x$ en $X$ y una sección regular $t\in\Gamma(U,\sh{O}_X)$ tal que $\sh{J}_X(D)|U=\sh{O}_U.t$ \sref{IV.21.2.7};
en otras palabras, la inmersión canónica $Y(D)\to X$ es \emph{regular y de codimensión $1$} en todo punto de $Y(D)$ \sref{IV.19.1.4}.
Recíprocamente, si $Y$ es un subpreesquema cerrado de $X$, inmerso regularmente en $X$ y de codimensión $1$ en todo punto de $Y$, existe un único divisor positivo $D$ tal que $Y(D)=Y$, pues para todo $x\in Y$ existe un entorno $U$ de $x$ en $X$ tal que $Y\cap U$ está definido por un ideal de $\sh{O}_U$ de la forma $\sh{O}_U.t$, donde $t$ es regular en $\Gamma(U,\sh{O}_X)$.
Obsérvese que entonces se tiene $\operatorname{Supp}(D)=Y(D)$, pues decir que $D_x\neq0$ significa que $t_x$ (con las notaciones anteriores) no es invertible, es decir, que $x\in Y(D)$.
\end{env}

\oldpage[IV-4]{263}

